\documentclass[./examen.tex]{subfiles}
\usepackage[first=1, last=20]{lcg}
\newcommand{\random}{\rand\arabic{rand}}
\usepackage{verbatim}
\newcounter{z}
\newcommand\y[2]{
  \loop \ifnum\value{z} < #1
    #2%
    \stepcounter{z}%
  \repeat
}

\begin{document}
\begin{comment}
    
    \question Responder las siguientes preguntas:
    \begin{parts}
    \part[1] ¿Cómo define el flujo magnético y cuáles son sus unidades? 
    \part[1] ¿Qué son las líneas de fuerzas?. Explique y dibuje.
    \part[1] ¿Qué son los imanes y que tipo hay?
    \part[1] ¿Qué es un campo magnético y como se dá en una espira?
    \end{parts}
       \end{comment}


%\info{Evaluación Electrotécnia II}{Campo magnético, Inducción magnética, permeabilidad magnética,imanes, líneas de campo magnético, flujo magnético, unidades magnéticas}
\y{31}{
\begin{center}

\makebox[\textwidth]{Nombre y Apellido:\enspace\hrulefill}
\end{center}
\begin{questions}
   
    \question Resolver los siguientes problemas:
     \begin{parts}
     \part Un solenoide de $l=\random0cm$ de largo y radio de $r=0,\random m$ está formado por $N=\random00$ espiras y lo recorre una intensidad de $I=\random A$.
        \begin{subparts}
        \subpart[3] Calcular la intensidad de campo magnético $H$ en el interior del solenoide.
        \subpart[3] Calcular la inducción magnética $B$ siendo un núcleo de aire $\mu_0=4 \pi x 10^{-7} T\cdot m /A$.
        \end{subparts}
    \part[4] Calcular el flujo magnético $\Phi$ de un solenoide de superficie $S=\random cm^2$ siendo la inducción magnética de $B=\random000Gs$.
    \end{parts}
\end{questions}

}
\end{document}
